%%%%%%%%%%%%%%%%%%%%%%%%%%%%%%%%%%%%%%%%%%%%%%%%%%%%%%%%%%%%%
%% Begin exercise %%
%%%%%%%%%%%%%%%%%%%%%%%%%%%%%%%%%%%%%%%%%%%%%%%%%%%%%%%%%%%%%
\ex{Analog and Digital Conversions}
% \newif%\ifshowsolutions
% \showsolutionsfalse % <-- change to \showsolutionstrue to show

%%%%%%%%%%%%%%%%%%%%%%%%%%%%%%%%%%%%%%%%%%%%%%%%%%%%%%%%%%%%%
%% Task 4.1 %%
%%%%%%%%%%%%%%%%%%%%%%%%%%%%%%%%%%%%%%%%%%%%%%%%%%%%%%%%%%%%%

\task{Digital-to-Analog Converter (DAC) Output Calculation}

A 7-bit DAC has the digital input word $1100111_2$ and a reference voltage of $8 \ \text{V}$.\\

Calculate:

\subtask{Analog Output Voltage}

%\ifshowsolutions
\begin{solutionblock}

\begin{itemize}
    \item Number of bits: $n = 7$
    \item Reference voltage: $V_{\text{ref}} = 8 \ \text{V}$
\end{itemize}

The given binary input is:
\[
1100111_2
\]

Converting to decimal:
\[
1100111_2 = 1\cdot2^6 + 1\cdot2^5 + 0\cdot2^4 + 0\cdot2^3 + 1\cdot2^2 + 1\cdot2^1 + 1\cdot2^0
\]
\[
= 64 + 32 + 0 + 0 + 4 + 2 + 1 = 103
\]

For an ideal DAC, the output voltage is given by:
\[
V_{\text{out}} = \frac{D}{2^n} \, V_{\text{ref}}
\]

Substituting values:
\[
V_{\text{out}} = \frac{103}{2^7} \times 8
\]
\[
V_{\text{out}} = \frac{103}{128} \times 8
\]
\[
V_{\text{out}} \approx 6.4375 \ \text{V}
\]

\[
\boxed{V_{\text{out}} \approx 6.44 \ \text{V}}
\]

\end{solutionblock}
%\fi

\subtask{Conversion Resolution Voltage}

%\ifshowsolutions
\begin{solutionblock}

The resolution (voltage step size) of a DAC is defined as:
\[
\Delta V = \frac{V_{\text{ref}}}{2^n}
\]

Substituting the given values:
\[
\Delta V = \frac{8}{2^7} = \frac{8}{128}
\]
\[
\Delta V = 0.0625 \ \text{V}
\]

\[
\boxed{\Delta V = 62.5 \ \text{mV}}
\]

\end{solutionblock}
%\fi

%%%%%%%%%%%%%%%%%%%%%%%%%%%%%%%%%%%%%%%%%%%%%%%%%%%%%%%%%%%%%
%% Task 4.2 %%
%%%%%%%%%%%%%%%%%%%%%%%%%%%%%%%%%%%%%%%%%%%%%%%%%%%%%%%%%%%%%

\task{Analog-to-Digital Converter (ADC) Quantization}

An 8-bit ADC has a reference voltage of $10 \ \text{V}$ and an analogue input voltage of
$6.875 \ \text{V}$.\\

Calculate:

\subtask{Digital Output Word}

%\ifshowsolutions
\begin{solutionblock}

\begin{itemize}
    \item Number of bits: $n = 8$
    \item Reference voltage: $V_{\text{ref}} = 10 \ \text{V}$
    \item Input voltage: $V_{\text{in}} = 6.875 \ \text{V}$
\end{itemize}

The resolution (LSB size) of the ADC is:
\[
\Delta V = \frac{V_{\text{ref}}}{2^n} = \frac{10}{2^8}
\]
\[
\Delta V = \frac{10}{256} = 0.0390 \ \text{V}
\]

The corresponding digital code is given by:
\[
D = \frac{V_{\text{in}}}{\Delta V}
\]

\[
D = \frac{6.875}{0.0390} = 176
\]

Converting the decimal value to binary:
\[
176_{10} = 10110000_2
\]

\[
\boxed{\text{Digital Output Word} = 10110000_2}
\]

\end{solutionblock}
%\fi

\subtask{Percentage Resolution}

%\ifshowsolutions
\begin{solutionblock}

\[
\text{Percentage Resolution} = \frac{1}{2^n} \times 100\% = \frac{1}{256} \times 100\%
\]

\[
\boxed{\text{Percentage Resolution} \approx 0.39\%}
\]

\end{solutionblock}
%\fi

%%%%%%%%%%%%%%%%%%%%%%%%%%%%%%%%%%%%%%%%%%%%%%%%%%%%%%%%%%%%%
%% Task 4.2 %%
%%%%%%%%%%%%%%%%%%%%%%%%%%%%%%%%%%%%%%%%%%%%%%%%%%%%%%%%%%%%%

\task{DAC Reference Voltage Calculation}

A 6-bit DAC is required to produce an output voltage of $9 \ \text{V}$ when all six
input bits are at logic 1.\\

Calculate the required reference voltage for this DAC.

\subtask{Reference Voltage}

%\ifshowsolutions
\begin{solutionblock}

\begin{itemize}
    \item Number of bits: $n = 6$
    \item Full-scale digital input: $D_{\text{max}} = 2^6 - 1 = 63$
    \item Required output voltage: $V_{\text{out, max}} = 9 \ \text{V}$
\end{itemize}

For an ideal DAC, the output voltage is given by:
\[
V_{\text{out}} = \frac{D}{2^n} \, V_{\text{ref}}
\]

At full-scale input ($D = D_{\text{max}} = 63$):
\[
9 = \frac{63}{64} \, V_{\text{ref}}
\]

Solving for the reference voltage:
\[
V_{\text{ref}} = \frac{9 \times 64}{63}
\]
\[
V_{\text{ref}} \approx 9.1429 \ \text{V}
\]

\[
\boxed{V_{\text{ref}} \approx 9.14 \ \text{V}}
\]

\end{solutionblock}
%\fi

%%%%%%%%%%%%%%%%%%%%%%%%%%%%%%%%%%%%%%%%%%%%%%%%%%%%%%%%%%%%%
%% Task 4.4 %%
%%%%%%%%%%%%%%%%%%%%%%%%%%%%%%%%%%%%%%%%%%%%%%%%%%%%%%%%%%%%%

\task{4-bit Weighted DAC Error Analysis}

A 4-bit weighted inverting DAC has the following parameters:

\begin{itemize}
    \item Reference voltage: $V_R = 1 \ \text{V}$
    \item Feedback resistor: $R_f = 8 R$
    \item Input resistors (MSB → LSB): $R$, $2R$, $4R$, $8R$
\end{itemize}

We calculate the output voltage $V_{\text{out}}$ step by step using the current summation method.

\subtask{Find currents through each resistor}
\begin{solutionblock}
Each bit contributes a current through its resistor if the bit is logic 1:

\[
I_n = \frac{V_R}{R_n}
\]

\begin{center}
\begin{tabular}{|c|c|c|}
\hline
Bit & Resistor & Current if bit = 1 (units of $V_R/R$) \\
\hline
$b_3$ (MSB) & $R$   & $I_1 = 1$ \\
$b_2$        & $2R$  & $I_2 = 0.5$ \\
$b_1$        & $4R$  & $I_3 = 0.25$ \\
$b_0$ (LSB) & $8R$  & $I_4 = 0.125$ \\
\hline
\end{tabular}
\end{center}
\end{solutionblock}

\subtask{Total current and output voltage}
\begin{solutionblock}

The total current entering the summing node:

\[
I_f = I_1 + I_2 + I_3 + I_4
\]

The inverting output voltage is:

\[
V_{\text{out}} = - R_f \, I_f
\]

Since $R_f = 8R$ and currents are in units of $V_R/R$, we have:

\[
V_{\text{out}} = -8 \times (I_1 + I_2 + I_3 + I_4) \ \text{V}
\]
\end{solutionblock}

\subtask{Compute $V_{\text{out}}$ for all 16 inputs}

\begin{solutionblock}
\begin{center}
\begin{tabular}{|c|c|c|}
\hline
Binary Input & Sum of Currents $I_f$ & $V_{\text{out}}$ (V) \\
\hline
0000 & 0     & 0   \\
0001 & 0.125 & -1  \\
0010 & 0.25  & -2  \\
0011 & 0.375 & -3  \\
0100 & 0.5   & -4  \\
0101 & 0.625 & -5  \\
0110 & 0.75  & -6  \\
0111 & 0.875 & -7  \\
1000 & 1     & -8  \\
1001 & 1.125 & -9  \\
1010 & 1.25  & -10 \\
1011 & 1.375 & -11 \\
1100 & 1.5   & -12 \\
1101 & 1.625 & -13 \\
1110 & 1.75  & -14 \\
1111 & 1.875 & -15 \\
\hline
\end{tabular}
\end{center}
\end{solutionblock}

\subtask{Compare with given DAC outputs}
\begin{solutionblock}

We now compare the calculated outputs ($V_{\text{out}}$) with the measured outputs:
\\
NOTE: We compare the absolute vaules of the actual outputs since the negative signs 
defines the nature of inverting amplifier DAC.
\begin{center}
\begin{tabular}{|c|c|c|c|}
\hline
\textbf{Binary Input} & \textbf{Calculated $V_{\text{out}}$ (V)} & \textbf{Given $V_{\text{out}}$ (V)} & \textbf{Comment} \\
\hline
0000 & 0   & 0   &  Correct \\
0001 & 1   & 1   &  Correct \\
0010 & 2   & 0   &  Incorrect \\
0011 & 3   & 1   &  Incorrect \\
0100 & 4   & 4   &  Correct \\
0101 & 5   & 5   &  Correct \\
0110 & 6   & 4   &  Incorrect \\
0111 & 7   & 5   &  Incorrect \\
1000 & 8   & 8   &  Correct \\
1001 & 9   & 9   &  Correct \\
1010 & 10  & 8   &  Incorrect \\
1011 & 11  & 9   &  Incorrect \\
1100 & 12  & 12  &  Correct \\
1101 & 13  & 13  &  Correct \\
1110 & 14  & 12  &  Incorrect \\
1111 & 15  & 13  &  Incorrect \\
\hline
\end{tabular}
\end{center}
\end{solutionblock}

\subtask{Assuming the electronic switch is working correctly, state the 
most likely cause of the DAC functioning incorrectly. }

\begin{solutionblock}
\begin{itemize}
    \item All incorrect outputs occur when the \textbf{second least significant bit} ($b_1$, weight $4R$) is logic 1.
    \item Since the electronic switches are assumed working correctly, the most likely cause is:
    \[
    \boxed{\text{Open-circuit or incorrect resistor in the $b_1$ branch (weight $4R$)}}.
    \]
    \item The DAC fails to add the $b_1$ contribution, producing the missing 2 V steps in the output.
\end{itemize}
\end{solutionblock}

%%%%%%%%%%%%%%%%%%%%%%%%%%%%%%%%%%%%%%%%%%%%%%%%%%%%%%%%%%%%%
%% Task 4.5 %%
%%%%%%%%%%%%%%%%%%%%%%%%%%%%%%%%%%%%%%%%%%%%%%%%%%%%%%%%%%%%%

\task{R-2R Ladder DAC Output Calculation}

A 4-bit R-2R ladder DAC has the following parameters:

\begin{itemize}
    \item Reference voltage: $V_{\text{ref}} = 10 \ \text{V}$
    \item Ladder resistor: $R = 20 \ \text{k}\Omega$
    \item Feedback resistor: $R_f = 50 \ \text{k}\Omega$
    \item Op-amp supply limits: $\pm 18 \ \text{V}$
\end{itemize}

\subtask{Output Voltage for Input 1001}

%\ifshowsolutions
\begin{solutionblock}

For an inverting R-2R ladder DAC, the output voltage is given by:

\[
V_{\text{out}} = - V_{\text{ref}} \frac{R_f}{R} \left( \frac{b_3}{2} + \frac{b_2}{4} + \frac{b_1}{8} + \frac{b_0}{16} \right)
\]

Where $b_3$ is the MSB and $b_0$ is the LSB.

The input word is:
\[
1001 \quad \Rightarrow \quad b_3=1, \ b_2=0, \ b_1=0, \ b_0=1
\]

\[
V_{\text{out}} = - -10 \cdot \frac{50\text{k}\Omega}{20\text{k}\Omega} \left( \frac{1}{2} + 0 + 0 + \frac{1}{16}  \right)
\]

\[
\boxed{V_{\text{out}} \approx -14.06 \ \text{V}}
\]

\end{solutionblock}
%\fi

\subtask{Maximum Digital Input Word That Can Be Converted}

%\ifshowsolutions
\begin{solutionblock}

The maximum output voltage of the DAC is limited by the op-amp supply: $|V_{\text{out}}| \le 18 \ \text{V}$.

The full-scale output (all bits = 1) is:

\[
V_{\text{FS}} = - V_{\text{ref}} \frac{R_f}{R} \left( \frac{1}{2} + \frac{1}{4} + \frac{1}{8} + \frac{1}{16} \right)
\]

\[
V_{\text{FS}} = -10 \cdot 2.5 \cdot \frac{15}{16} = -25 \cdot 0.9375
\]

\[
V_{\text{FS}} = -23.4375 \ \text{V}
\]

Since $|V_{\text{FS}}| = 23.44 \ \text{V} > 18 \ \text{V}$, the op-amp saturates.  

To find the maximum digital word that does not saturate, let $D_{\text{max}}$ be the decimal equivalent of the digital input:

\[
V_{\text{out}} = - V_{\text{ref}} \frac{R_f}{R} \frac{D_{\text{max}}}{2^N} \le 18
\]

\[
25 \cdot \frac{D_{\text{max}}}{16} \le 18
\]

\[
D_{\text{max}} \le \frac{18 \cdot 16}{25} = 11.52
\]

\[
\Rightarrow D_{\text{max}} = 11 \quad (\text{round down to nearest integer})
\]

Convert to 4-bit binary:

\[
D_{\text{max}} = 1011
\]

\[
\boxed{\text{Maximum input word that can be converted without saturation: } 1011}
\]

\end{solutionblock}
%\fi


%%%%%%%%%%%%%%%%%%%%%%%%%%%%%%%%%%%%%%%%%%%%%%%%%%%%%%%%%%%%%
%% Task 4.6 %%
%%%%%%%%%%%%%%%%%%%%%%%%%%%%%%%%%%%%%%%%%%%%%%%%%%%%%%%%%%%%%

\task{AD Conversion with successive approximation}

A 12-bit converter with an input voltage range of 0-10V converts
a voltage of 7.1875V.

\subtask{Calculate the number of conversion steps (Clock cycles) 
necessary to get the exact result as a digital number.}

%\ifshowsolutions
\begin{solutionblock}

This converter needs a start bit. If the bit is set exactly one 
conversion is performed. In the case of the input voltage chosen, 
the DA converter output is exactly matching with it after 4 bits 
have been tested. Every additional voltage out of the DA converter
is too much. That means the remaining bits will be all "0". The 
converter always tests all bits. That means the conversion doesn't
stop if a match appear !\\
The conversion time is always the same.\\
For 12-bit resolution, it needs 12 cycles regardless of the input value.

\[
V_{\text{out}} = V_{\text{ref}} \left(
\frac{b_{N-1}}{2} + \frac{b_{N-2}}{4} + \cdots + \frac{b_1}{2^{N-1}} + \frac{b_0}{2^N}
\right)
\]

\[
V_{\text{out}} = V_{\text{ref}} \left(
2^{N-1} b_{N-1} + 2^{N-2} b_{N-2} + \cdots + 2 b_1 + b_0
\right) / 2^N
\]

\[
V_{\text{out}} = V_{\text{ref}}
D
/ 2^N
\]

The resolution of an $N$-bit ADC is:

\[
\Delta V = \frac{V_{\text{FS}}}{2^N - 1}
\]

where $V_{\text{FS}}$ is the full-scale voltage.  

\[
\Delta V = \frac{10}{2^{12}-1} = \frac{10}{4095} \approx 0.002442 \ \text{V per step}
\]

The digital output code corresponding to $V_{\text{in}}$ is:

\[
D = \frac{V_{\text{in}}}{\Delta V}
\]

\[
D = \frac{7.1875}{0.002442} \approx 2943
\]

For a successive approximation ADC, the number of steps to reach this output is equal to the digital code plus 1 (starting from zero):

\[
\text{Clock cycles} = D + 1 = 2943 + 1 = 2944
\]

\begin{figure}[h!]
    \centering
    \includegraphics[width=1\textwidth]{successive_approximation.pdf}
    \caption{Successive Approximation ADC}
\end{figure}

\end{solutionblock}